\makeatletter
\@ifundefined{keepfig}{\newenvironment{keepfig}{\par\edef\keepfig@pi{\the\parindent}\noindent
  \begin{minipage}{\linewidth}\parindent=\keepfig@pi\relax}{\end{minipage}\par}}{}
\makeatother
\section{Симметрические многочлены}

Пусть $R$~--- коммутативное кольцо с единицей. Многочлен $P(X_1,\dots,X_n)\in R[X_1,\dots,X_n]$ называется
\emph{симметрическим}, если он не меняется при перестановках переменных:
\[
P(X_{\sigma(1)},\dots,X_{\sigma(n)})=P(X_1,\dots,X_n)\qquad\text{для всех }\sigma\in S_n .
\]
Например, симметрическими являются многочлены
\begin{align*}
p_1&=\textstyle\sum_{i}X_i=X_1+X_2+\dots+X_n,\\
p_2&=\textstyle\sum_{i<j}X_iX_j=X_1X_2+X_1X_3+\dots+X_1X_n+X_2X_3+\dots+X_{n-1}X_n,\\
p_3&=\textstyle\sum_{i<j<k}X_iX_jX_k=X_1X_2X_3+\cdots,\\
&\ \,\vdots\\
p_r&=\textstyle\sum_{i_1<\dots<i_r}X_{i_1}\cdots X_{i_r},\\
&\ \,\vdots\\
p_n&=X_1X_2\cdots X_n:
\end{align*}
$p_r$~--- сумма \emph{всех} мономов степени $r$, являющихся произведениями $r$ различных переменных, а перестановка
переменных лишь переставляет эти мономы. Многочлены $p_1,\dots,p_n$ называются \emph{элементарными симметрическими
многочленами}. Раскрывая скобки, получаем тождество в $R[X_1,\dots,X_n][T]$:
\begin{equation*}
(T-X_1)(T-X_2)\cdots(T-X_n)=T^n-p_1T^{n-1}+p_2T^{n-2}-\dots+(-1)^np_n
\end{equation*}
(коэффициент при $T^{n-r}$ получается выбором слагаемого $-X_i$ ровно в $r$ скобках).

\subsection*{Основная теорема}

\begin{theorem}[5.36: основная теорема о симметрических многочленах]
Каждый симметрический многочлен $P\in R[X_1,\dots,X_n]$ является многочленом от элементарных симметрических
многочленов с коэффициентами из $R$, т.\,е.\ $P\in R[p_1,\dots,p_n]$.
\end{theorem}

Идея доказательства: упорядочим мономы и будем <<съедать>> старший моном многочлена $P$, вычитая подходящее произведение
степеней $p_i$ с тем же старшим мономом. Старший моном при этом строго уменьшается, а бесконечно уменьшаться он не может.

\paragraph{Порядок на мономах.} Положим
\[
X_1^{i_1}X_2^{i_2}\cdots X_n^{i_n}>X_1^{j_1}X_2^{j_2}\cdots X_n^{j_n},
\]
если либо $i_1+\dots+i_n>j_1+\dots+j_n$, либо суммы равны и для некоторого $s$
\[
(i_1,\dots,i_s)=(j_1,\dots,j_s),\qquad\text{но}\qquad i_{s+1}>j_{s+1}
\]
(сначала сравниваем полные степени, а при равных степенях~--- показатели лексикографически). Например,
\[
X_1X_2X_3^3>X_1X_2^2X_3>X_1X_2X_3^2 .
\]
Это линейный порядок на множестве мономов со следующими свойствами.
\begin{enumerate}[label=\textup{(\roman*)}]
\item \emph{Он согласован с умножением:} если $m>m'$, то $mu>m'u$ для любого монома $u$ (умножение на $u$ прибавляет к
обоим наборам показателей один и тот же набор, а это не меняет ни сравнения полных степеней, ни лексикографического
сравнения).
\item \emph{Ниже каждого монома $m$ лежит лишь конечное число мономов:} все они имеют полную степень $\le\deg m$. Поэтому
каждое непустое множество мономов имеет наименьший элемент.
\end{enumerate}
Для ненулевого $P$ \emph{старшим мономом} назовём наибольший моном $m$, входящий в $P$ с ненулевым коэффициентом $c$,
а $cm$~--- \emph{старшим членом}. Из (i)
следует: если $P$ и $Q$ имеют старшие члены $cm$ и $c'm'$ и $cc'\ne0$, то старший член $PQ$ равен $cc'mm'$. Действительно,
произведение любых других двух мономов $u\le m$, $v\le m'$ (хотя бы одно неравенство строгое) меньше $mm'$: например, при
$u<m$ имеем $uv<mv\le mm'$; поэтому член $cc'mm'$ ни с чем не сокращается. В частности, это так, если оба старших
коэффициента равны $1$ (делители нуля в $R$ здесь не мешают).

\begin{proof}[Доказательство теоремы \textup{5.36}]
\emph{Шаг 1: старший моном симметрического многочлена.} Пусть $P\ne0$ симметричен и
\[
P=cX_1^{i_1}\cdots X_n^{i_n}+\text{(младшие члены)},\qquad c\ne0 .
\]
Так как $P$ симметричен, вместе с мономом $X_1^{i_1}\cdots X_n^{i_n}$ он содержит (с тем же коэффициентом $c$) все мономы,
получающиеся из него перестановками переменных. Если бы $i_k<i_{k+1}$ для некоторого $k$, то транспозиция $X_k\leftrightarrow
X_{k+1}$ дала бы моном той же степени, у которого первые $k-1$ показателей те же, а $k$-й больше,~--- т.\,е.\ моном
\emph{выше} старшего. Значит,
\[
i_1\ge i_2\ge\dots\ge i_n .
\]

\emph{Шаг 2: старший моном произведения элементарных многочленов.} Старший моном $p_k$~--- это $X_1X_2\cdots X_k$, с
коэффициентом $1$. По свойству (i) старший член произведения $p_1^{d_1}p_2^{d_2}\cdots p_n^{d_n}$ равен
\begin{equation}\label{eq:s11-leadmon}
X_1^{d_1+d_2+\dots+d_n}\,X_2^{d_2+\dots+d_n}\cdots X_n^{d_n}
\end{equation}
(с коэффициентом $1$): переменная $X_k$ входит в старшие мономы ровно тех множителей $p_r$, у которых $r\ge k$.

\emph{Шаг 3: уменьшение старшего монома.} Положим $d_k=i_k-i_{k+1}$ ($k<n$) и $d_n=i_n$; по шагу 1 все $d_k\ge0$. Тогда
$d_k+d_{k+1}+\dots+d_n=i_k$, и по \eqref{eq:s11-leadmon} многочлен $c\,p_1^{i_1-i_2}p_2^{i_2-i_3}\cdots p_n^{i_n}$ имеет тот
же старший член $cX_1^{i_1}\cdots X_n^{i_n}$, что и $P$, а все его остальные мономы ниже. Поэтому многочлен
\begin{equation}\label{eq:s11-reduce}
P(X_1,\dots,X_n)-c\,p_1^{i_1-i_2}p_2^{i_2-i_3}\cdots p_n^{i_n}
\end{equation}
симметричен (разность симметрических) и либо равен нулю, либо имеет старший моном строго ниже, чем у $P$.

\emph{Шаг 4: завершение.} Предположим, что теорема неверна. По свойству (ii) среди симметрических многочленов, не лежащих в
$R[p_1,\dots,p_n]$, найдётся многочлен $P$ с наименьшим старшим мономом. Многочлен \eqref{eq:s11-reduce} симметричен и
имеет меньший старший моном (или равен нулю), поэтому лежит в $R[p_1,\dots,p_n]$. Но тогда и $P$ лежит в
$R[p_1,\dots,p_n]$~--- противоречие.
\end{proof}

\begin{remark}[5.37]
(a) \emph{Доказательство даёт алгоритм.} Возьмём, например,
\begin{align*}
P(X_1,X_2)&=(X_1+7X_1X_2+X_2)^2\\
&=X_1^2+2X_1X_2+14X_1^2X_2+X_2^2+14X_1X_2^2+49X_1^2X_2^2 .
\end{align*}
Старший член~--- $49X_1^2X_2^2$ (показатели $(2,2)$), поэтому вычитаем $49p_2^2$:
\[
P-49p_2^2=X_1^2+2X_1X_2+14X_1^2X_2+X_2^2+14X_1X_2^2 .
\]
Теперь старший член $14X_1^2X_2$ (показатели $(2,1)$), вычитаем $14p_1^{2-1}p_2^{1}=14p_1p_2$:
\[
P-49p_2^2-14p_1p_2=X_1^2+2X_1X_2+X_2^2,
\]
и наконец $P-49p_2^2-14p_1p_2-p_1^2=0$. Итак, $P=p_1^2+14p_1p_2+49p_2^2$ (что, конечно, видно и сразу:
$P=(p_1+7p_2)^2$).

Ещё пример: $P=X_1^3+X_2^3+X_3^3$ при $n=3$. Старший моном $X_1^3$, показатели $(3,0,0)$, вычитаем $p_1^3$:
\[
P-p_1^3=-3\sum_{i\ne j}X_i^2X_j-6X_1X_2X_3 .
\]
Старший член теперь $-3X_1^2X_2$, показатели $(2,1,0)$; прибавляем $3p_1p_2$, где
\[
p_1p_2=\sum_{i\ne j}X_i^2X_j+3X_1X_2X_3 ,
\]
и получаем
\[
P-p_1^3+3p_1p_2=3X_1X_2X_3=3p_3 .
\]
Итак, $X_1^3+X_2^3+X_3^3=p_1^3-3p_1p_2+3p_3$.

(b) \emph{Представление $P$ в виде многочлена от $p_i$ единственно.} Иначе, вычитая, мы получили бы ненулевой многочлен
$Q(t_1,\dots,t_n)=\sum_dc_d\,t_1^{d_1}\cdots t_n^{d_n}$ с $Q(p_1,\dots,p_n)=0$ в $R[X_1,\dots,X_n]$. Но старшие мономы
\eqref{eq:s11-leadmon} произведений $p_1^{d_1}\cdots p_n^{d_n}$ попарно различны: отображение
$(d_1,\dots,d_n)\mapsto(d_1+\dots+d_n,\ d_2+\dots+d_n,\ \dots,\ d_n)$ инъективно ($d_k$ восстанавливается как разность
соседних координат). Возьмём среди наборов $d$ с $c_d\ne0$ тот, у которого моном \eqref{eq:s11-leadmon} наибольший. Этот
моном входит в $c_d\,p_1^{d_1}\cdots p_n^{d_n}$ с коэффициентом $c_d\ne0$, а во все остальные слагаемые $Q(p_1,\dots,p_n)$
не входит (их мономы не выше их старших мономов, а те меньше выбранного). Значит, $Q(p_1,\dots,p_n)\ne0$~--- противоречие.
Иначе говоря, \emph{элементарные симметрические многочлены алгебраически независимы над $R$}; это понадобится в
\rf{5.40}.

(c) \emph{Симметрические функции от корней.} Пусть
\[
f(X)=X^n+a_1X^{n-1}+\dots+a_n\in R[X]
\]
и $f$ раскладывается на линейные множители над некоторым кольцом $S\supset R$:
$f(X)=\prod_{i=1}^n(X-\alpha_i)$, $\alpha_i\in S$. Сравнивая с тождеством для $\prod(T-X_i)$ выше, получаем
\[
a_1=-p_1(\alpha_1,\dots,\alpha_n),\quad a_2=p_2(\alpha_1,\dots,\alpha_n),\quad\dots,\quad
a_n=(-1)^np_n(\alpha_1,\dots,\alpha_n).
\]
Итак, элементарные симметрические функции от корней $f$ лежат в $R$, и по теореме \rf{5.36} для любого симметрического
$P\in R[X_1,\dots,X_n]$ значение $P(\alpha_1,\dots,\alpha_n)$ лежит в $R$: если $P=Q(p_1,\dots,p_n)$, то
$P(\alpha_1,\dots,\alpha_n)=Q(-a_1,a_2,\dots,(-1)^na_n)$. Например, \emph{дискриминант}
\[
D(f)\overset{\mathrm{def}}{=}\prod_{i<j}(\alpha_i-\alpha_j)^2
\]
лежит в $R$; более того, $D(f)=\Delta_n(a_1,\dots,a_n)$ для многочлена $\Delta_n$ с целыми коэффициентами, не зависящего
от $R$ (применяем теорему при $R=\Z$ к симметрическому многочлену $\prod_{i<j}(X_i-X_j)^2\in\Z[X_1,\dots,X_n]$). При $n=2$
из $(X_1-X_2)^2=p_1^2-4p_2$ получаем привычное $D(X^2+bX+c)=b^2-4c$.
\end{remark}

\subsection*{Симметрические рациональные функции}

\begin{theorem}[5.38: теорема о симметрических функциях]
Пусть $F$~--- поле, и $S_n$ действует на поле рациональных функций $F(X_1,\dots,X_n)$ перестановками переменных. Тогда
поле инвариантов равно $F(p_1,\dots,p_n)$:
\[
F(X_1,\dots,X_n)^{S_n}=F(p_1,\dots,p_n).
\]
\end{theorem}

\begin{proof}
Включение $\supset$ очевидно: $p_i$ симметричны. Обратно, пусть $f\in F(X_1,\dots,X_n)$ симметрична (неподвижна
относительно $S_n$). Запишем $f=g/h$, $g,h\in F[X_1,\dots,X_n]$, $h\ne0$. Многочлен
\[
H=\prod_{\sigma\in S_n}\sigma h
\]
симметричен: перестановка $\tau$ лишь переставляет множители $\sigma h\mapsto(\tau\sigma)h$. Произведение
$Hf=g\cdot\prod_{\sigma\ne1}\sigma h$~--- тоже многочлен, и он симметричен как произведение симметрических элементов поля.
По теореме \rf{5.36} оба многочлена $H$ и $Hf$ лежат в $F[p_1,\dots,p_n]$, поэтому $f=Hf/H\in F(p_1,\dots,p_n)$.
\end{proof}

\begin{corollary}[5.39]
Поле $F(X_1,\dots,X_n)$ является расширением Галуа поля $F(p_1,\dots,p_n)$ с группой Галуа $S_n$ (действующей
перестановками переменных $X_i$).
\end{corollary}

\begin{proof}
Перестановки переменных задают инъективный гомоморфизм $S_n\to\Aut(F(X_1,\dots,X_n))$ (разные перестановки по-разному
двигают переменные), так что $S_n$~--- конечная группа автоморфизмов поля $F(X_1,\dots,X_n)$. По теореме \rf{5.38}
$F(p_1,\dots,p_n)=F(X_1,\dots,X_n)^{S_n}$, и по теореме Артина (\rf{3.11}) $F(X_1,\dots,X_n)$~--- расширение Галуа поля
$F(p_1,\dots,p_n)$ с группой $S_n$; в частности, $[F(X_1,\dots,X_n):F(p_1,\dots,p_n)]=n!$.
\end{proof}

Поле $F(X_1,\dots,X_n)$ является полем разложения над $F(p_1,\dots,p_n)$ многочлена
\[
g(T)=(T-X_1)\cdots(T-X_n)=T^n-p_1T^{n-1}+\dots+(-1)^np_n\in F(p_1,\dots,p_n)[T]:
\]
$g$ раскладывается в $F(X_1,\dots,X_n)$ на линейные множители, и его корни $X_1,\dots,X_n$ порождают это поле (даже над
$F$). Поэтому группа Галуа многочлена $g$ над $F(p_1,\dots,p_n)$~--- это $S_n$, действующая на его корнях
$X_1,\dots,X_n$ всеми возможными перестановками.

\begin{remark}[5.39a: второе доказательство 5.38 и 5.39]
Теорему \rf{5.38} (а вместе с ней и \rf{5.39}) можно получить, не пользуясь теоремой \rf{5.36},~--- одним подсчётом степеней. Рассмотрим
башню
\[
F(p_1,\dots,p_n)\subset F(X_1,\dots,X_n)^{S_n}\subset F(X_1,\dots,X_n).
\]
Верхнее поле~--- поле разложения многочлена $g$ степени $n$ над нижним, поэтому
\[
[F(X_1,\dots,X_n):F(p_1,\dots,p_n)]\le n!
\]
(присоединение $X_k$ к $F(p_1,\dots,p_n)(X_1,\dots,X_{k-1})$ даёт степень $\le n-k+1$, так как $X_k$~--- корень
многочлена $g/\bigl((T-X_1)\cdots(T-X_{k-1})\bigr)$ степени $n-k+1$ с коэффициентами в этом поле; степени
перемножаются по \rf{1.20}, как в доказательстве \rf{2.7}). С другой стороны, по теореме Артина (\rf{3.11})
$[F(X_1,\dots,X_n):F(X_1,\dots,X_n)^{S_n}]=|S_n|=n!$. По мультипликативности
степеней (\rf{1.20}) $[F(X_1,\dots,X_n)^{S_n}:F(p_1,\dots,p_n)]\le1$, т.\,е.\ $F(X_1,\dots,X_n)^{S_n}=F(p_1,\dots,p_n)$.
\medskip
\begin{keepfig}
\begin{center}
\begin{tikzpicture}
\node (A) at (0,0) {$F(p_1,\dots,p_n)$};
\node at (1.55,0) {$\subset$};
\node (B) at (3.4,0) {$F(X_1,\dots,X_n)^{S_n}$};
\node at (5.25,0) {$\subset$};
\node (C) at (6.9,0) {$F(X_1,\dots,X_n)$};
\draw[decorate,decoration={brace,mirror,amplitude=5pt},cnew] (2.1,-0.35) -- node[below=6pt,font=\footnotesize]{$n!$ (теорема Артина)} (8.1,-0.35);
\draw[decorate,decoration={brace,amplitude=5pt},cold] (-1.1,0.35) -- node[above=6pt,font=\footnotesize]{$\le n!$ (поле разложения многочлена степени $n$)} (8.1,0.35);
\end{tikzpicture}
\end{center}
\figcap{Сверху степень всей башни не больше $n!$, снизу степень верхнего этажа ровно $n!$; поэтому нижний этаж
тривиален.}
\end{keepfig}
Тот же приём~--- <<чтобы доказать, что поле мало, найди много автоморфизмов, оставляющих его на месте, и сравни степени>>~---
часто применяется вместе с теоремой Артина.
\end{remark}

\paragraph{Исторический комментарий.} Симметрические многочлены играли важную роль в работах Галуа. В его
мемуаре об условиях разрешимости уравнений в радикалах доказывается следующее утверждение. \emph{Пусть $f$~---
многочлен с коэффициентами $\sigma_1,\dots,\sigma_n$, $x_1,\dots,x_n$~--- его корни, а $U,V,\dots$~--- некоторые
рациональные функции от $x_i$. Тогда существует группа $G$ перестановок корней $x_i$ такая, что рациональные функции
от $x_i$, неподвижные относительно всех перестановок из $G$,~--- это в точности функции, рационально выражающиеся через
$\sigma_1,\dots,\sigma_n$ и $U,V,\dots$} Если в качестве $U,V,\dots$ взять элементы промежуточного поля $E$ между полем
коэффициентов $f$ и полем разложения $f$, то утверждение говорит, что существует группа $G$ перестановок корней, неподвижное
поле которой (при действии на поле разложения) есть в точности $E$,~--- это одна из половин основной теоремы теории
Галуа (\rf{3.17}).
